\chapter{Synchronization}

\section{Race conditions and critical sections}

Threads in a process share global data. Without coordination, concurrent access gives results
that depend on timing.

\begin{example}[Lost update]
Two threads each run \code{total++} 1000 times on a shared \code{int total = 0}. The statement
compiles to roughly three instructions:
\begin{lstlisting}[style=mips]
lw   t0, total      # load
addi t0, t0, 1      # increment
sw   t0, total      # store
\end{lstlisting}
Interleaving: T1 loads 0, T2 loads 0, T1 stores 1, T2 stores 1. One increment is lost. The final
value can be anything from 2 to 2000 (2 needs a tricky interleaving, but is possible).
\end{example}

\begin{definition}
A \textbf{race condition} occurs when the result of a program depends on the relative timing
or interleaving of threads. A \textbf{critical section} is a piece of code that accesses shared
data and must not be executed by more than one thread at a time. \textbf{Mutual exclusion} is
the property that at most one thread is inside the critical section.
\end{definition}

Requirements for a correct critical section solution:
\begin{enumerate}
  \item \textbf{Mutual exclusion}: at most one thread inside.
  \item \textbf{Progress}: if no thread is inside and some want in, one of them eventually
    gets in.
  \item \textbf{Bounded waiting} (fairness): a thread that wants in is not overtaken forever.
\end{enumerate}

\begin{aside}
\textbf{Compiler and CPU reordering.} The compiler may keep a shared variable in a register or
reorder memory accesses; \code{volatile} forces every access to go to memory but does
\emph{not} make operations atomic. Modern CPUs also reorder loads/stores; synchronization
primitives contain \emph{memory barriers} (\code{sync} on MIPS) so that accesses do not leak
out of critical sections.
\end{aside}

\section{Achieving mutual exclusion}

\subsection{Disabling interrupts}
On a \emph{single} processor, a thread can only lose the CPU through an interrupt. Disabling
interrupts (\code{splhigh()} in OS/161) around a critical section therefore gives mutual
exclusion. Drawbacks: only the kernel may do it; it does not work on multiprocessors; it delays
device handling, so critical sections must be very short.

\subsection{Peterson's algorithm (software only)}
For two threads, using only loads and stores (assuming sequential consistency):
\begin{lstlisting}[style=c]
volatile bool flag[2] = {false, false};
volatile int turn;

void enter(int i) {          /* i is 0 or 1 */
    int j = 1 - i;
    flag[i] = true;          /* I want in */
    turn = j;                /* but you go first */
    while (flag[j] && turn == j) { /* busy wait */ }
}
void leave(int i) { flag[i] = false; }
\end{lstlisting}
Correct but impractical: two threads only, busy waits, and breaks under memory reordering.

\subsection{Hardware atomic instructions}
\begin{definition}[Test-and-set]
\code{TestAndSet(addr, v)} atomically stores \code{v} at \code{addr} and returns the old value.
\end{definition}
x86 provides \code{xchg}; many architectures provide \code{compare-and-swap}. MIPS instead has a
pair: \code{ll} (load-linked) and \code{sc} (store-conditional). \code{sc} only succeeds if no
other write to that address happened since the matching \code{ll}; it returns success/failure.

\begin{lstlisting}[style=c]
/* test-and-set built from ll/sc, as in OS/161's spinlock_data_testandset */
static inline int test_and_set(volatile int *addr) {
    int x, y;
    y = 1;
    __asm volatile(
        ".set push; .set mips32;"
        "ll %0, 0(%2);"       /* x = *addr (linked) */
        "sc %1, 0(%2);"       /* try *addr = 1; y = success? */
        ".set pop"
        : "=&r"(x), "+r"(y) : "r"(addr));
    if (y == 0) return 1;     /* sc failed: pretend lock was held */
    return x;                 /* old value */
}
\end{lstlisting}

\section{Spinlocks}

\begin{lstlisting}[style=c]
void spinlock_acquire(struct spinlock *lk) {
    splraise(IPL_NONE, IPL_HIGH);      /* disable interrupts on this CPU */
    for (;;) {
        if (lk->lk_lock != 0) continue;            /* test first (cheap read) */
        if (test_and_set(&lk->lk_lock) == 0) break; /* got it */
    }
    lk->lk_holder = curcpu->c_self;
    membar_any_any();
}
void spinlock_release(struct spinlock *lk) {
    membar_any_any();
    lk->lk_holder = NULL;
    lk->lk_lock = 0;
    spllower(IPL_HIGH, IPL_NONE);
}
\end{lstlisting}

\begin{itemize}
  \item A spinlock \emph{busy-waits}: the waiting thread burns CPU.
  \item OS/161 spinlocks also disable interrupts on the local CPU, so the holder cannot be
    preempted (which would make others spin for a whole quantum) and an interrupt handler on the
    same CPU cannot deadlock trying to take the same lock.
  \item Use spinlocks only for very short critical sections. Never block (sleep) while holding
    a spinlock.
\end{itemize}

\section{Locks (mutexes)}

A \textbf{lock} provides mutual exclusion like a spinlock, but a thread that cannot acquire it
\emph{blocks} instead of spinning. Locks also have an \textbf{owner}: only the thread that
acquired a lock may release it.

\begin{lstlisting}[style=c]
struct lock {
    char *lk_name;
    struct wchan *lk_wchan;
    struct spinlock lk_spin;   /* protects the fields below */
    struct thread *lk_owner;   /* NULL if free */
};

void lock_acquire(struct lock *lk) {
    spinlock_acquire(&lk->lk_spin);
    while (lk->lk_owner != NULL) {
        /* atomically releases lk_spin and sleeps; reacquires on wakeup */
        wchan_sleep(lk->lk_wchan, &lk->lk_spin);
    }
    lk->lk_owner = curthread;
    spinlock_release(&lk->lk_spin);
}

void lock_release(struct lock *lk) {
    KASSERT(lock_do_i_hold(lk));
    spinlock_acquire(&lk->lk_spin);
    lk->lk_owner = NULL;
    wchan_wakeone(lk->lk_wchan, &lk->lk_spin);
    spinlock_release(&lk->lk_spin);
}
\end{lstlisting}

\begin{important}[Why \code{while} and not \code{if}?]
After being woken, another thread may grab the lock before the woken thread runs. The woken
thread must re-check the condition. Also, \code{wchan\_sleep} must release the spinlock and
sleep \emph{atomically}, otherwise a wakeup could happen between the release and the sleep and
be lost.
\end{important}

\begin{center}\small
\begin{tabular}{@{}lll@{}}
\toprule
 & Spinlock & Lock \\
\midrule
Waiting & busy-wait & block on wait channel \\
Good for & very short sections, interrupt handlers & longer sections \\
Can sleep while holding? & no & yes \\
Overhead & low when uncontended & context switch when contended \\
\bottomrule
\end{tabular}
\end{center}

\section{Semaphores}

\begin{definition}[Semaphore (Dijkstra)]
A \textbf{semaphore} is an object with a non-negative integer value and two atomic operations:
\begin{itemize}
  \item \code{P} (wait, ``proberen''): if the value is 0, block until it is positive; then
    decrement it.
  \item \code{V} (signal, ``verhogen''): increment the value (waking a waiter if any).
\end{itemize}
\end{definition}

\begin{lstlisting}[style=c]
void P(struct semaphore *sem) {
    spinlock_acquire(&sem->sem_lock);
    while (sem->sem_count == 0)
        wchan_sleep(sem->sem_wchan, &sem->sem_lock);
    sem->sem_count--;
    spinlock_release(&sem->sem_lock);
}
void V(struct semaphore *sem) {
    spinlock_acquire(&sem->sem_lock);
    sem->sem_count++;
    wchan_wakeone(sem->sem_wchan, &sem->sem_lock);
    spinlock_release(&sem->sem_lock);
}
\end{lstlisting}

Kinds of use:
\begin{itemize}
  \item \textbf{Binary semaphore} (initial value 1): like a lock, but no owner --- any thread
    can \code{V}.
  \item \textbf{Counting semaphore} (initial value $N$): $N$ identical resources.
  \item \textbf{Signalling} (initial value 0): one thread waits (\code{P}) for an event that
    another thread announces (\code{V}). E.g.\ a parent waits for $k$ children by calling
    \code{P} $k$ times; each child calls \code{V} once when done.
\end{itemize}

\subsection{Producer/consumer (bounded buffer)}
\begin{lstlisting}[style=c]
#define N 16
item buffer[N]; int in = 0, out = 0;
struct semaphore *items;   /* initial 0: number of full slots */
struct semaphore *spaces;  /* initial N: number of empty slots */
struct lock *mutex;        /* protects buffer, in, out */

void produce(item x) {
    P(spaces);                    /* wait for an empty slot */
    lock_acquire(mutex);
    buffer[in] = x; in = (in + 1) % N;
    lock_release(mutex);
    V(items);                     /* announce a full slot */
}
item consume(void) {
    P(items);
    lock_acquire(mutex);
    item x = buffer[out]; out = (out + 1) % N;
    lock_release(mutex);
    V(spaces);
    return x;
}
\end{lstlisting}
\begin{remark}
Swapping the order (taking \code{mutex} before \code{P(spaces)}) can deadlock: a producer holding
the mutex sleeps on a full buffer, and no consumer can get the mutex to free a slot.
\end{remark}

\section{Condition variables}

A \textbf{condition variable} lets a thread wait, \emph{inside} a critical section protected by a
lock, until some condition on shared data becomes true. It is always used together with a lock.

\begin{description}
  \item[\code{cv\_wait(cv, lk)}] atomically release \code{lk} and block on \code{cv}; when woken,
    reacquire \code{lk} before returning.
  \item[\code{cv\_signal(cv, lk)}] wake one thread waiting on \code{cv} (caller holds \code{lk}).
  \item[\code{cv\_broadcast(cv, lk)}] wake all threads waiting on \code{cv}.
\end{description}

\begin{lstlisting}[style=c]
void cv_wait(struct cv *cv, struct lock *lk) {
    spinlock_acquire(&cv->cv_spin);
    lock_release(lk);                         /* release only after holding cv_spin */
    wchan_sleep(cv->cv_wchan, &cv->cv_spin);  /* so no signal can be lost */
    spinlock_release(&cv->cv_spin);
    lock_acquire(lk);
}
\end{lstlisting}

OS/161 (like POSIX) uses \textbf{Mesa semantics}: a signalled thread is merely made ready; by the
time it runs the condition may be false again. Hence the golden rule:

\begin{important}[Always wait in a loop]
\begin{lstlisting}[style=c]
lock_acquire(lk);
while (!condition)        /* NOT "if" */
    cv_wait(cv, lk);
/* condition holds and we hold lk */
lock_release(lk);
\end{lstlisting}
(Hoare semantics, where the signaller hands the lock directly to the waiter, would allow
\code{if}, but almost no real system implements it.)
\end{important}

\begin{example}[Bounded buffer with a lock and two CVs]
\begin{lstlisting}[style=c]
int count = 0;
struct lock *mtx; struct cv *notfull, *notempty;

void produce(item x) {
    lock_acquire(mtx);
    while (count == N) cv_wait(notfull, mtx);
    buffer[in] = x; in = (in + 1) % N; count++;
    cv_signal(notempty, mtx);
    lock_release(mtx);
}
item consume(void) {
    lock_acquire(mtx);
    while (count == 0) cv_wait(notempty, mtx);
    item x = buffer[out]; out = (out + 1) % N; count--;
    cv_signal(notfull, mtx);
    lock_release(mtx);
    return x;
}
\end{lstlisting}
\end{example}

\section{Deadlock}

\begin{example}
Thread 1: \code{lock\_acquire(A); lock\_acquire(B);} \quad
Thread 2: \code{lock\_acquire(B); lock\_acquire(A);}. If each gets its first lock, both wait
forever.
\end{example}

\begin{definition}[Deadlock]
A set of threads is \textbf{deadlocked} if each is blocked waiting for a resource held by
another thread in the set, so none can make progress.
\end{definition}

\begin{theorem}[Coffman conditions]
Deadlock can occur only if all four hold simultaneously:
\begin{enumerate}
  \item \textbf{Mutual exclusion} (no sharing): a resource can be held by only one thread.
  \item \textbf{Hold and wait}: a thread holds resources while requesting more.
  \item \textbf{No preemption}: resources cannot be forcibly taken away.
  \item \textbf{Circular wait}: a cycle of threads each waiting for the next.
\end{enumerate}
\end{theorem}

\subsection{Prevention} Break one condition:
\begin{itemize}
  \item \emph{No hold and wait}: request all resources at once, or release everything before
    requesting more (wasteful, may starve).
  \item \emph{Preemption}: if a request cannot be satisfied, release what you hold and retry
    (e.g.\ \code{lock\_tryacquire} then back off). Risk of livelock.
  \item \emph{Resource ordering}: number all resources and always acquire in increasing
    order. Then a cycle is impossible: in a cycle some thread would hold a higher-numbered
    resource while waiting for a lower one. This is the most common practical technique.
\end{itemize}

\subsection{Detection and recovery}
Model the system as a \textbf{resource allocation graph}: an edge $T \to R$ if $T$ requests $R$,
$R \to T$ if $R$ is allocated to $T$. With single-instance resources, deadlock $\iff$ a cycle
exists. With multiple instances, use a detection algorithm:

\begin{lstlisting}[style=c]
/* U[i][j]: units of resource j requested by thread i
   A[i][j]: units of resource j allocated to thread i
   avail[j]: free units of resource j */
bool done[n] = {false};
bool progress = true;
while (progress) {
    progress = false;
    for (i = 0; i < n; i++)
        if (!done[i] && U[i] <= avail) {   /* componentwise compare */
            avail += A[i];                 /* pretend T_i finishes, frees all */
            done[i] = true; progress = true;
        }
}
/* every thread with done[i] == false is deadlocked */
\end{lstlisting}

\begin{example}
Resources $R_1,R_2$ with 1 and 2 units. $A = \{T_1:(1,0), T_2:(0,1), T_3:(0,1)\}$, so
avail $=(0,0)$. Requests $U = \{T_1:(0,1), T_2:(1,0), T_3:(0,0)\}$. $T_3$ can finish (needs
nothing): avail $=(0,1)$. Now $T_1$ can finish: avail $=(1,1)$. Then $T_2$. No deadlock.
If instead $T_3$ also requested $(1,0)$, nobody could proceed: $T_1,T_2,T_3$ deadlocked.
\end{example}

Recovery: kill a thread in the cycle, or preempt and roll back. Many systems simply ignore
deadlock (the ``ostrich algorithm'') and rely on lock-ordering discipline in code.

\section*{Summary}
\begin{itemize}
  \item Races arise when interleavings of accesses to shared data produce different results;
    critical sections must be mutually exclusive.
  \item Hardware atomics (test-and-set, CAS, \code{ll}/\code{sc}) build spinlocks; spinlocks
    plus wait channels build blocking locks, semaphores and condition variables.
  \item Semaphores: \code{P}/\code{V} on a counter; good for counting resources and signalling.
  \item CVs: wait for a condition while holding a lock; Mesa semantics $\Rightarrow$ wait in a
    \code{while} loop.
  \item Deadlock needs mutual exclusion, hold-and-wait, no preemption, circular wait; prevent by
    lock ordering, or detect with a graph / matrix algorithm.
\end{itemize}
